\section{Opposite labels on actual belt-cell sides}
\label{sec:area_law_belt_fan_colors}

Use the selected opposing region to assign labels to each belt-cell fan,
as in \cite[Section~11]{OpenAI2026AreaLaw}. The residue choices at every
layer are arbitrary; the label assignment requires no sparsity estimate.

% Source: 10-geometry.tex, prop:two-families, lines 172--174,
% 212--218 and 299--323, revision adc7f1241b42e322a6451854ab7e4b4c146bf78a.
% These are local assignment results, not the full two-family proposition.

\begin{definition}[Labels of a belt-cell fan]
    \label{def:area_law_belt_fan_color}
    \lean{TNLean.PEPS.AreaLaw.Geometry.beltCellFanColor}
    \leanok
    \uses{def:area_law_selected_elementary_opponent,def:area_law_belt_cells,
        def:area_law_actual_side_mask}
    Fix a common origin $o\in\mathbb R^2$, a finite endpoint set $Z$,
    integer width $C_0\ge2$ and natural lower index $k_0\ge50\,000\,000$.
    For every layer $h$, choose residues $0\le a_h,b_h<m_h$, where
    $m_h=2^{p_h-\ell_h}$, $\ell_h$ is the fine exponent and $p_h$ is the
    pitch exponent. Let $\mathcal B_h$ be the corresponding actual belt-cell
    index set. For $k\ge k_0$ and $z\in\mathcal B_k$, denote the elementary
    segments of the actual fan by $e_i=[a_i,b_i]$ and their selected opposing
    identifiers by $P_{k,z}(i)$. Order the triples $(k,z_1,z_2)$
    lexicographically. Define $c_{k,z}(i)\in\mathbb Z/2\mathbb Z$ by
    \begin{align}
        c_{k,z}(i)=
        \begin{cases}
            (k_0-1)+1, & P_{k,z}(i)=\star,                                              \\
            h+1,       & P_{k,z}(i)=(h,w),\ w\notin\mathcal B_h,                        \\
            0,         & P_{k,z}(i)=(h,w),\ w\in\mathcal B_h,\ (k,z_1,z_2)<(h,w_1,w_2), \\
            1,         & P_{k,z}(i)=(h,w),\ w\in\mathcal B_h,\ (k,z_1,z_2)>(h,w_1,w_2).
        \end{cases}
    \end{align}
    The indexed cells are distinct whenever the opponent is a fine cell,
    so the last two alternatives are exhaustive.
\end{definition}

\begin{theorem}[The label opposite the dummy tile]
    \label{thm:area_law_belt_fan_dummy_color}
    \lean{TNLean.PEPS.AreaLaw.Geometry.beltCellFanColor_dummy}
    \leanok
    \uses{def:area_law_belt_fan_color,def:area_law_dyadic_layers}
    Retain the data of the label definition. If an actual elementary segment
    $e_i$ lies in $\overline{N_{k_0}}$, then
    $c_{k,z}(i)=(k_0-1)+1\ne k_0-1$ in $\mathbb Z/2\mathbb Z$.
    Thus its incident triangle has the label opposite the dummy tile.
\end{theorem}
\begin{proof}\leanok
    \uses{thm:area_law_unique_elementary_side_opponent,def:area_law_belt_fan_color}
    The dummy identifier satisfies the whole-segment containment condition.
    Uniqueness therefore gives $P_{k,z}(i)=\star$. The asserted label follows
    from its definition; adding one changes each residue modulo two.
\end{proof}

\begin{theorem}[The label and retained primary of a non-belt opponent]
    \label{thm:area_law_belt_fan_nonbelt_color}
    \lean{TNLean.PEPS.AreaLaw.Geometry.beltCellFanColor_nonbelt_opponent}
    \leanok
    \uses{def:area_law_belt_fan_color,def:area_law_nonbelt_pitch_index,
        def:area_law_primary_pitch_indices,def:area_law_primary_regions}
    Retain the data of the label definition. Let $h\ge k_0$ and let
    $w\in F_{h,\ell_h}\setminus\mathcal B_h$ be an actual non-belt index.
    If $e_i\cap\overline{Q_{\ell_h}(w)}$ contains two distinct points,
    put $J=j(w)$, the shifted pitch index at residues $a_h,b_h$.
    Then $J$ is retained in the actual primary index set $\mathcal P_h$,
    the whole segment $e_i$ lies in the primary birth region
    $\overline{D_h\cap I_{a_h,b_h}(J)}$, and
    $c_{k,z}(i)=h+1\ne h$ in $\mathbb Z/2\mathbb Z$.
    The primary identifier is $(h,J)$, including when its region is disconnected.
\end{theorem}
\begin{proof}\leanok
    \uses{thm:area_law_selected_opponent_contact,
        thm:area_law_unique_elementary_side_opponent,
        thm:area_law_nonbelt_primary_containment,
        thm:area_law_primary_pitch_membership,def:area_law_belt_fan_color}
    Belt and non-belt membership imply $(k,z)\ne(h,w)$. The contact
    characterization gives $P_{k,z}(i)=(h,w)$, and the selected opponent
    contains the whole segment in its closed square. Closed non-belt-cell
    containment consequently gives
    $e_i\subseteq\overline{Q_{\ell_h}(w)}\subseteq
        \overline{D_h\cap I_{a_h,b_h}(J)}$.
    Since $e_i$ contains its first endpoint, this last closure is nonempty.
    The layer--pitch intersection itself is therefore nonempty, and the exact
    retained-index characterization gives $J\in\mathcal P_h$.
    The non-belt clause in the label definition gives $h+1$, which differs from $h$.
\end{proof}

\begin{theorem}[Opposite labels across contacting belt-cell segments]
    \label{thm:area_law_belt_fan_contact_colors}
    \lean{TNLean.PEPS.AreaLaw.Geometry.beltCellFanColor_belt_contact}
    \leanok
    \uses{def:area_law_belt_fan_color}
    Retain the common data and residue choices. Let $(k,z)\ne(h,w)$
    be actual belt-cell pairs at $k,h\ge k_0$, and let $e_i,e_j$
    be arbitrary elementary segments in their respective actual subdivisions.
    If $e_i\cap e_j$ contains two distinct points, then
    $c_{k,z}(i)\ne c_{h,w}(j)$.
    No matching slot or reversed endpoints are assumed.
\end{theorem}
\begin{proof}\leanok
    \uses{thm:area_law_selected_opponent_contact,thm:area_law_cell_fan_cover,
        def:area_law_belt_fan_color}
    Each elementary segment lies in its incident triangle: the segment's convex
    hull is contained in the convex hull of the three triangle vertices.
    The fan cover then places it in its own closed cell.
    Thus the assumed intersection gives two-point contact of each segment
    with the other closed cell. The contact characterization yields
    $P_{k,z}(i)=(h,w)$ and $P_{h,w}(j)=(k,z)$.
    The distinct lexicographic keys have opposite strict orders, so the two
    belt clauses assign zero and one in opposite order.
\end{proof}

The local labels can be used in the equal-label runs of each fan.
Global primary--run interfaces, point contacts, boundary sectors and the
isolated-star construction remain further parts of the two-family construction.
